\begin{papercard}{vonelm2007}
\cf{Stages} \stagetag{R} (reporting items for a descriptive, non-experimental study), \stagetag{A} (setting, eligibility and selection of units), \stagetag{D} (variables, data sources, measurement, missing data), \stagetag{S} (study design stated early). The STROBE Statement (\emph{PLoS Medicine} 4(10):e296, 2007; 5 PDF pages including references): a 22-item checklist for reporting cohort, case-control and cross-sectional studies. It is for epidemiology; the thesis' sample of apps and documents is not a study of persons, so the card lists which items transfer to a descriptive, cross-sectional study of apps and where the guideline does not fit. The companion Explanation and Elaboration paper is not in the file.

\cf{Problem and aims} Much biomedical research is observational, yet reporting is often inadequate, which ``hampers the assessment of its strengths and weaknesses and of a study's generalisability'' \pc{vonelm2007}{1}{Much biomedical research is observational}. Evidence of gaps cited: confounder choice often not reported; and in a survey of longitudinal stroke studies 17 of 49 articles (35\,\%) did not specify eligibility criteria \pc{vonelm2007}{2}{In a survey of longitudinal studies in stroke research, 17 of 49 articles (35\%) did not specify the eligibility criteria}. The intention is ``solely to provide guidance on how to report observational research well'' and the checklist ``is not an instrument to evaluate the quality'' of such research \pc{vonelm2007}{2}{The intention is solely to provide guidance on how to report observational research well}; \pc{vonelm2007}{2}{the checklist is not an instrument to evaluate the quality of observational research}.

\cf{Development} STROBE was set up in 2004; the group searched textbooks, databases, reference lists and personal files, and chose the three designs used most widely in analytical observational research \pc{vonelm2007}{2}{We decided to focus on the 3 study designs that are used most widely in analytical observational research}. A 2-day workshop in Bristol (September 2004) had 23 attendees (journal editors, epidemiologists, methodologists, statisticians, practitioners from Europe and North America) \pc{vonelm2007}{2}{We organised a 2-day workshop in Bristol, UK, in September 2004. 23 individuals attended this meeting}; three working groups drafted items per design, a combined draft went online, three revisions and two comment summaries were published, and the coordinating group met eight times for one or two days \pc{vonelm2007}{2}{met on eight occasions for 1 or 2 days}. No formal Delphi or voting statistics are reported.

\cf{The 22 items (Table 1)} Items 1 (title/abstract), 2--3 (introduction), 4--12 (methods), 13--17 (results), 18--21 (discussion), 22 (funding) \pc{vonelm2007}{2}{The STROBE Statement is a checklist of 22 items that we consider essential for good reporting of observational studies}; 18 items are common to all designs and four (6, 12, 14, 15) have design-specific versions \pc{vonelm2007}{4}{while four (items 6, 12, 14, and 15)}.
\begin{itemize}\itemsep0pt
\item \emph{1} state the design in title or abstract \pc{vonelm2007}{3}{Indicate the study’s design with a commonly used term in the title or the abstract} and give an informative, balanced abstract. \emph{2--3} background and rationale; objectives ``including any prespecified hypotheses'' \pc{vonelm2007}{3}{including any prespecified hypotheses}.
\item \emph{4} key design elements early; \emph{5} setting, locations and relevant dates (recruitment, exposure, follow-up, data collection); \emph{6} eligibility criteria, sources and methods of selection (for cross-sectional studies: eligibility, sources and selection methods); \emph{7} define all outcomes, exposures, predictors, confounders and effect modifiers \pc{vonelm2007}{3}{Clearly define all outcomes, exposures, predictors, potential confounders, and effect modifiers}; \emph{8} data sources and measurement methods for each variable \pc{vonelm2007}{3}{For each variable of interest, give sources of data and details of methods of assessment}; \emph{9} efforts to address bias \pc{vonelm2007}{3}{Describe any efforts to address potential sources of bias}; \emph{10} how the study size was arrived at \pc{vonelm2007}{3}{Explain how the study size was arrived at}; \emph{11} how quantitative variables were handled and grouped; \emph{12} statistical methods including subgroups, missing data \pc{vonelm2007}{3}{Explain how missing data were addressed} and sensitivity analyses.
\item \emph{13} numbers at each stage, reasons for non-participation, optional flow diagram \pc{vonelm2007}{3}{Report the numbers of individuals at each stage of the study}; \emph{14} characteristics and number with missing data per variable \pc{vonelm2007}{3}{Indicate the number of participants with missing data for each variable of interest}; \emph{15} outcome data; \emph{16} unadjusted and adjusted estimates with precision; \emph{17} other analyses.
\item \emph{18} key results; \emph{19} limitations, ``both direction and magnitude of any potential bias'' \pc{vonelm2007}{3}{Discuss both direction and magnitude of any potential bias}; \emph{20} cautious interpretation considering multiplicity of analyses and other evidence \pc{vonelm2007}{3}{Give a cautious overall interpretation of results considering objectives, limitations, multiplicity of analyses}; \emph{21} generalisability (external validity) \pc{vonelm2007}{3}{Discuss the generalisability (external validity) of the study results}; \emph{22} funding and the funders' role \pc{vonelm2007}{3}{Give the source of funding and the role of the funders for the present study}.
\end{itemize}

\cf{Intended use and non-prescriptive stance} The statement helps authors, editors, reviewers and critical readers; it covers exploratory studies, hypothesis-driven analyses of available data and planned data collection, because readers ``always need to know what was planned (and what was not), what was done, what was found, and what the results mean'' \pc{vonelm2007}{4}{Some studies are essentially exploratory and raise interesting hypotheses}. Items can be addressed anywhere in the article and in any order; the authors reject standardising text: ``manuscripts should not be STROBEd'' \pc{vonelm2007}{4}{We believe that manuscripts should not be STROBEd, in the sense of regulating style or terminology}, and they encourage narrative elements. They also state the checklist is not a tool for assessing quality and that examples in the companion paper were chosen for reporting, not for the quality of the research \pc{vonelm2007}{4}{We emphasise that the STROBE Statement was not developed as a tool for assessing the quality of published observational research}.

\cf{Authors' stated limits} Restricted to three designs (extensions invited, e.g.\ for case-crossover or ecological studies; STREGA for genetics is underway); no systematic review was done for each item and no research to fill evidence gaps \pc{vonelm2007}{4}{We did not undertake a comprehensive systematic review for each of the checklist items and sub-items}; contributors came through existing networks and were dominated by Europe and North America, and probably not representative of disciplines \pc{vonelm2007}{4}{was not representative in terms of geography (it was dominated by contributors from Europe and North America)}; the checklist is an evolving document. Better reporting may help temper over-enthusiastic reporting of new findings and support decisions about whether new studies are needed.

\cf{Limits (card)} Written for studies of people, exposures and outcomes: confounding, effect estimates, follow-up and matching items (6b, 12d, 16) do not apply to a descriptive analysis of apps and documents. There is no evidence in the paper that using the checklist improves reporting (it points to CONSORT's effects only by citation). The explanatory paper that defines each item is separate and not in the corpus, so wording such as ``participants'' must be translated by the user. A guideline for \emph{reporting} cannot say how coding of policies or apps should be done.

\cf{Use for this thesis} (1) The empirical part (a sample of health apps and their privacy documents, coded and counted) is best described as a cross-sectional descriptive study of apps; STROBE items can structure the methods and results sections where PRISMA-ScR (for literature) does not fit: item 5 setting and dates (stores, country, capture dates), item 6 eligibility and selection, items 7 and 8 variable definitions and measurement (the codebook and who or what coded), item 9 bias, item 10 sample size and how it was reached, item 12(c) and 14(b) missing data and number of apps with missing fields, item 13 flow of apps from identified to included with reasons, item 19 limitations with direction of bias, item 21 generalisability (the sample is a selection from stores, not all apps). (2) Treat ``participants'' as ``apps'' and ``documents'' and state this mapping. (3) The warning against ``STROBEing'' a manuscript supports adapting items instead of reproducing a checklist. (4) Combine with PRISMA-S for the search log and O'Brien et al.\ (SRQR) for the qualitative coding parts, and say which items from each were used.
\end{papercard}
